\section{Quick Setup Guide}
This project isolates dependencies using a Python virtual environment to ensure absolute reproducibility. 

\subsection{Environment Setup}
To initialize the project on a new machine, execute the following commands in the root of the project:

\begin{lstlisting}[language=bash]
# 1. Create a Python Virtual Environment
|\BeginAccSupp{ActualText={\detokenize{py -3.14 -m venv .venv}}}\ttfamily\footnotesize py -3.14 -m venv .venv\EndAccSupp{}|

# 2. Activate the Environment (Windows PowerShell)
|\BeginAccSupp{ActualText={\detokenize{.\.venv\Scripts\Activate.ps1}}}\ttfamily\footnotesize .\textbackslash .venv\textbackslash Scripts\textbackslash Activate.ps1\EndAccSupp{}|

# 3. Install all dependencies
|\BeginAccSupp{ActualText={\detokenize{pip install -r requirements.txt}}}\ttfamily\footnotesize pip install -r requirements.txt\EndAccSupp{}|
\end{lstlisting}

\subsection{Dependencies}
The \texttt{requirements.txt} installs the following libraries:
\begin{itemize}
    \item \texttt{numpy} \& \texttt{pandas}: Fast numerical array manipulation and CSV template handling.
    \item \texttt{torch} (PyTorch): Deep learning framework for the 1D-CNN. Installed from the PyTorch Nightly channel with CUDA 12.8 support via \texttt{-{}-extra-index-url}.
    \item \texttt{plotly} \& \texttt{kaleido}: Interactive HTML visualizations and static PNG export.
    \item \texttt{scikit-learn}: Utility functions for metrics and data splitting.
    \item \texttt{matplotlib}: Supplementary plotting library.
    \item \texttt{scipy}, \texttt{statsmodels}, \texttt{tslearn}: Statistical interpolation and legacy dataset compatibility.
\end{itemize}

\noindent\textbf{Note:} The \texttt{-{}-pre} flag in \texttt{requirements.txt} enables installation of pre-release PyTorch builds, which is required for CUDA 12.8 support on newer GPUs (e.g., NVIDIA RTX 5060).
